\chapter{\MakeUppercase{Conclusão}}

Este trabalho apresentou o \textbf{SIM --- Sistema de Informações Municipais}, plataforma conversacional multi-instância para acesso a dados de transparência pública, com \textbf{instância piloto em Caeté (MG)} integrando recuperação documental sobre o Diário Oficial e consulta analítica via Text-to-SQL sobre bases SICOM/CGU.

A arquitetura híbrida --- orquestração por LLM, RAG e DuckDB --- demonstra viabilidade técnica de reduzir a barreira entre o cidadão e dados municipais abertos. A validação com 80 consultas de referência comprovou geração confiável de SQL (VSR de 97,5\%) e equivalência funcional (EF) em 62,5\% dos casos no baseline; o refinamento do \textit{prompt} (v2.1) elevou a EF para 72,5\% com VSR integral. A validação piloto do RAG ($N_{\text{RAG}} = 8$) obteve Recall@5 de 100\% e MRR de 63,33\% na busca densa; melhorias de recuperação (deduplicação, busca híbrida por número de ato e filtro por tipo) elevaram o MRR a 87,5\% no mesmo \textit{golden}.

As contribuições incluem: (i)~conceito SIM multi-instância e arquitetura replicável para GovTech municipal; (ii)~\textit{pipeline} ETL e exposição analítica sobre dados SICOM/CGU; (iii)~protocolo de validação reprodutível versionado em \texttt{avaliacao/}; (iv)~evidências empíricas para escolha de modelo e tamanho de contexto; (v)~validação complementar da recuperação documental; e (vi)~repositório público para replicação.

Como trabalhos futuros, propõe-se: ampliação do \textit{golden} RAG e julgamento de respostas em linguagem natural; revisão humana do conjunto \textit{golden} SQL; testes com usuários reais; replicação do SIM em outros municípios; e autocorreção iterativa de consultas SQL.
